The measuring instrument is a persona subspace: four directions in the
base model's layer-22 activations that capture persona-like variation,
i.e.\ whether text reads as the assistant or as some other character
speaking. One direction is the assistant-contrast axis of
\citet{lu2026assistant}; the other three are the top principal components
of activations across the 275-role inventory from
\citet{song2026role}, reused verbatim and hash-pinned. I extracted the
activations and constructed this four-direction basis for the present study.
The layer was selected by a pre-registered d-prime criterion
(d-prime: a standard measure of how cleanly two score distributions
separate) inside the L12--L22 band. Each response or document is represented
by a 4,096-number mean-pooled
activation vector. ``Frozen'' means the whole specification was committed
before the relevant computation ran, so it could not be tuned to an
outcome. Two derived quantities carry the analysis:
\emph{persona load}, the fraction of a direction's norm that lies inside
the persona subspace ($0$ = none of it, $1$ = all of it), and
\emph{separability}, the residual outside it
($\sqrt{1 - \mathrm{load}^2}$; $1.0$ = fully persona-orthogonal). The
persona-separability hypothesis supplied the pre-registered H2 order; the
measured index values were not used to write it. This run does not validate
those values as a direct or out-of-sample filterability ranker.
